%%%%%%%%%%%%%%%%%%%%%%%%%%%%%%%%%%%%%%%%%%%%%%%%%%%%%%%%%%%%%%%%%%%%%%%%%%%%%%%
%% TKS FORMAL MATHEMATICAL MANUAL v7.4 — BEAUTIFUL CORRECTED EDITION
%%
%% CANONICAL INHERITANCE:
%%   v6.1 -> v7.0 -> v7.1 -> v7.2 -> v7.3 -> v7.4
%%
%% FEATURES:
%%   - Beautiful colored theorem environments throughout
%%   - Charter font for elegant typography
%%   - All mathematical corrections integrated
%%   - Professional formatting with no text overflow
%%%%%%%%%%%%%%%%%%%%%%%%%%%%%%%%%%%%%%%%%%%%%%%%%%%%%%%%%%%%%%%%%%%%%%%%%%%%%%%

\documentclass[11pt,twoside,openright]{book}

%% ============================================================================
%% PACKAGES
%% ============================================================================
\usepackage[utf8]{inputenc}
\usepackage[T1]{fontenc}
% Use default Computer Modern font (widely available)
\usepackage{amsmath,amssymb,amsthm,amsfonts}
\usepackage{mathtools}
%\usepackage{stmaryrd}                   % For semantic brackets
\usepackage{bm}                         % Bold math
\usepackage{enumitem}
\usepackage{booktabs}
\usepackage{array}
\usepackage{longtable}
\usepackage{multirow}
\usepackage{graphicx}
\usepackage{tikz}
\usetikzlibrary{cd,arrows,positioning,calc,decorations.pathmorphing,shapes}
\usepackage{hyperref}
\usepackage{cleveref}
\usepackage{xcolor}
\usepackage{tcolorbox}
\tcbuselibrary{theorems,skins,breakable}
\usepackage{fancyhdr}
\usepackage{titlesec}
\usepackage{geometry}
\usepackage{listings}
\usepackage{multicol}
\usepackage{braket}                     % For quantum notation
\usepackage{tikz-cd}                    % For commutative diagrams
%\usepackage{bussproofs}                 % For proof trees
\usepackage{mathrsfs}                   % For \mathscr (Scott topology)
%\usepackage{wasysym}                    % Additional symbols
%\usepackage{extarrows}                  % Extended arrows
%\usepackage{proof}                      % Proof derivations
%\usepackage{tensor}                     % For tensor notation
\usepackage{float}                      % For [H] float placement
\usepackage{adjustbox}                  % For scaling wide tables
\usepackage{tocloft}                    % For TOC formatting

%% ============================================================================
%% PAGE GEOMETRY AND OVERFLOW PREVENTION
%% ============================================================================
\geometry{
  paper=letterpaper,
  inner=1.2in,
  outer=1in,
  top=1in,
  bottom=1in
}

% Prevent text overflow
\setlength{\emergencystretch}{3em}
\tolerance=9999
\hbadness=9999
\sloppy

% Paragraph spacing
\setlength{\parskip}{0.6\baselineskip}

%% ============================================================================
%% COLORS
%% ============================================================================
\definecolor{defblue}{RGB}{30,90,150}
\definecolor{defbluebg}{RGB}{235,245,255}
\definecolor{thmgreen}{RGB}{34,139,34}
\definecolor{thmgreenbg}{RGB}{235,255,235}
\definecolor{lemgreen}{RGB}{60,179,113}
\definecolor{lemgreenbg}{RGB}{240,255,240}
\definecolor{propteal}{RGB}{0,128,128}
\definecolor{proptealbg}{RGB}{235,255,255}
\definecolor{axpurple}{RGB}{128,0,128}
\definecolor{axpurplebg}{RGB}{250,235,255}
\definecolor{remorange}{RGB}{255,140,0}
\definecolor{remorangebg}{RGB}{255,250,235}
\definecolor{expink}{RGB}{219,112,147}
\definecolor{expinkbg}{RGB}{255,240,245}
\definecolor{chapterblue}{RGB}{0,51,102}

%% ============================================================================
%% COLORED THEOREM ENVIRONMENTS
%% ============================================================================

% Counter for all theorem-like environments
\newcounter{thmcounter}[chapter]
\renewcommand{\thethmcounter}{\thechapter.\arabic{thmcounter}}

% Standard amsthm environments for included files
\theoremstyle{definition}
\newtheorem{definition}{Definition}[chapter]
\newtheorem{axiom}[definition]{Axiom}
\newtheorem{assumption}[definition]{Assumption}
\newtheorem{construction}[definition]{Construction}
\newtheorem{notation}[definition]{Notation}

\theoremstyle{plain}
\newtheorem{theorem}[definition]{Theorem}
\newtheorem{lemma}[definition]{Lemma}
\newtheorem{proposition}[definition]{Proposition}
\newtheorem{corollary}[definition]{Corollary}

\theoremstyle{remark}
\newtheorem{remark}[definition]{Remark}
\newtheorem{example}[definition]{Example}

% Definition - Blue (starred version for main document)
\newtcbtheorem[use counter=thmcounter]{definitionbox}{Definition}{
    enhanced,
    breakable,
    colback=defbluebg,
    colframe=defblue,
    fonttitle=\bfseries,
    left=8pt, right=8pt, top=6pt, bottom=6pt,
    boxrule=1pt,
    arc=2pt,
    separator sign={.},
}{def}

% Theorem - Green (starred version for main document)
\newtcbtheorem[use counter=thmcounter]{theorembox}{Theorem}{
    enhanced,
    breakable,
    colback=thmgreenbg,
    colframe=thmgreen,
    fonttitle=\bfseries,
    left=8pt, right=8pt, top=6pt, bottom=6pt,
    boxrule=1pt,
    arc=2pt,
    separator sign={.},
}{thm}

% Lemma - Light Green (starred version for main document)
\newtcbtheorem[use counter=thmcounter]{lemmabox}{Lemma}{
    enhanced,
    breakable,
    colback=lemgreenbg,
    colframe=lemgreen,
    fonttitle=\bfseries,
    left=8pt, right=8pt, top=6pt, bottom=6pt,
    boxrule=1pt,
    arc=2pt,
    separator sign={.},
}{lem}

% Proposition - Teal (starred version for main document)
\newtcbtheorem[use counter=thmcounter]{propositionbox}{Proposition}{
    enhanced,
    breakable,
    colback=proptealbg,
    colframe=propteal,
    fonttitle=\bfseries,
    left=8pt, right=8pt, top=6pt, bottom=6pt,
    boxrule=1pt,
    arc=2pt,
    separator sign={.},
}{prop}

% Corollary - Teal (starred version for main document)
\newtcbtheorem[use counter=thmcounter]{corollarybox}{Corollary}{
    enhanced,
    breakable,
    colback=proptealbg,
    colframe=propteal,
    fonttitle=\bfseries,
    left=8pt, right=8pt, top=6pt, bottom=6pt,
    boxrule=1pt,
    arc=2pt,
    separator sign={.},
}{cor}

% Axiom - Purple (starred version for main document)
\newtcbtheorem[use counter=thmcounter]{axiombox}{Axiom}{
    enhanced,
    breakable,
    colback=axpurplebg,
    colframe=axpurple,
    fonttitle=\bfseries,
    left=8pt, right=8pt, top=6pt, bottom=6pt,
    boxrule=1pt,
    arc=2pt,
    separator sign={.},
}{ax}

% Remark - Orange (starred version for main document)
\newtcbtheorem[use counter=thmcounter]{remarkbox}{Remark}{
    enhanced,
    breakable,
    colback=remorangebg,
    colframe=remorange,
    fonttitle=\bfseries,
    left=8pt, right=8pt, top=6pt, bottom=6pt,
    boxrule=1pt,
    arc=2pt,
    separator sign={.},
}{rem}

% Example - Pink (starred version for main document)
\newtcbtheorem[use counter=thmcounter]{examplebox}{Example}{
    enhanced,
    breakable,
    colback=expinkbg,
    colframe=expink,
    fonttitle=\bfseries,
    left=8pt, right=8pt, top=6pt, bottom=6pt,
    boxrule=1pt,
    arc=2pt,
    separator sign={.},
}{ex}

% Construction - Blue (starred version for main document)
\newtcbtheorem[use counter=thmcounter]{constructionbox}{Construction}{
    enhanced,
    breakable,
    colback=defbluebg,
    colframe=defblue,
    fonttitle=\bfseries,
    left=8pt, right=8pt, top=6pt, bottom=6pt,
    boxrule=1pt,
    arc=2pt,
    separator sign={.},
}{constr}

% Notation - Blue (starred version for main document)
\newtcbtheorem[use counter=thmcounter]{notationbox}{Notation}{
    enhanced,
    breakable,
    colback=defbluebg,
    colframe=defblue,
    fonttitle=\bfseries,
    left=8pt, right=8pt, top=6pt, bottom=6pt,
    boxrule=1pt,
    arc=2pt,
    separator sign={.},
}{not}

%% ============================================================================
%% CHAPTER AND SECTION FORMATTING
%% ============================================================================
\titleformat{\chapter}[display]
  {\normalfont\huge\bfseries\color{chapterblue}}
  {\chaptertitlename\ \thechapter}
  {20pt}
  {\Huge}
\titlespacing*{\chapter}{0pt}{50pt}{40pt}

\titleformat{\section}
  {\normalfont\Large\bfseries\color{chapterblue}}
  {\thesection}
  {1em}
  {}
  
\titleformat{\subsection}
  {\normalfont\large\bfseries\color{chapterblue!80}}
  {\thesubsection}
  {1em}
  {}

%% ============================================================================
%% TABLE OF CONTENTS FORMATTING
%% ============================================================================
\setlength{\cftbeforepartskip}{15pt}
\setlength{\cftbeforechapskip}{8pt}
\setlength{\cftbeforesecskip}{3pt}

%% ============================================================================
%% CUSTOM COMMANDS (TKS-specific)
%% ============================================================================
% Semantic brackets
\newcommand{\sem}[1]{[\![#1]\!]}

% TKS-specific symbols
\newcommand{\Worlds}{\mathcal{W}}
\newcommand{\Ideas}{\mathcal{I}}
\newcommand{\Elements}{\mathcal{E}}
\newcommand{\Noetics}{\mathcal{N}}
\newcommand{\Foundations}{\mathcal{F}}
\newcommand{\Acquisitions}{\mathcal{A}}
\newcommand{\Noetica}{\mathbf{Noetica}}
\newcommand{\Acq}{\mathbf{Acq}}

% Noetic operators
\newcommand{\noe}[1]{\nu_{#1}}

% Tootra operations
\newcommand{\Tadd}{\oplus_T}
\newcommand{\Tsub}{\ominus_T}
\newcommand{\Tmul}{\otimes_T}
\newcommand{\Tdiv}{\oslash_T}

% World association
\newcommand{\Wadd}{\oplus_W}
\newcommand{\Wsub}{\ominus_W}

% Type judgments
\newcommand{\types}{\vdash}
\newcommand{\typmark}{:}

% RPM Monad
\newcommand{\RPM}{\mathfrak{R}}
\newcommand{\rpmret}{\mathsf{return}}
\newcommand{\rpmbind}{\mathbin{\gg\!=}}
\newcommand{\Kleisli}{\mathcal{K}}

% Fractal notation
\newcommand{\Frac}{\Phi}
\newcommand{\FracSet}{\mathfrak{F}}
\newcommand{\bigstep}{\Downarrow}
\newcommand{\smallstep}{\rightarrow}
\newcommand{\Failed}{\mathsf{Failed}}

% Quantum notation
\newcommand{\Hspace}{\mathcal{H}}
\newcommand{\qket}[1]{|#1\rangle}
\newcommand{\qbra}[1]{\langle #1|}
\newcommand{\qbraket}[2]{\langle #1 | #2 \rangle}
\newcommand{\qop}[1]{\hat{#1}}

% Domain theory
\newcommand{\Scott}{\mathscr{S}}
\newcommand{\Dcpo}{\mathbf{Dcpo}}
\newcommand{\Cont}{\mathbf{Cont}}
\newcommand{\sqleq}{\sqsubseteq}
\newcommand{\fix}{\mathsf{fix}}

% Natural transformations
\newcommand{\NatTrans}[2]{\eta_{#1 \to #2}}
\newcommand{\WorldFunctor}[1]{\mathfrak{W}_{#1}}

% Transfinite fractals
\newcommand{\Ord}{\mathbf{Ord}}
\newcommand{\TransFrac}[1]{\Phi^{#1}}
\newcommand{\limfrac}{\varinjlim}
\newcommand{\colfrac}{\varprojlim}
\newcommand{\omegalim}{\omega\text{-}\lim}

% Topos and coalgebra
\newcommand{\Topos}{\mathbf{Topos}}
\newcommand{\Sh}{\mathbf{Sh}}
\newcommand{\Coalg}{\mathbf{Coalg}}
\newcommand{\Terminal}{\mathbf{1}}
\newcommand{\Initial}{\mathbf{0}}

% v7.4 New commands
\newcommand{\vnewfour}{\textbf{v7.4 New}}
\newcommand{\NoeticManifold}{M_\nu}
\newcommand{\TangentBundle}{TM_\nu}
\newcommand{\CotangentBundle}{T^*M_\nu}
\newcommand{\NoeticMetric}{g_\nu}
\newcommand{\NoeticConnection}{\nabla^\nu}
\newcommand{\NoeticCurvature}{R^\nu}
\newcommand{\InftyTopos}{\mathbf{Topos}_\infty}

% Info boxes
\newcommand{\infobox}[2]{%
\begin{tcolorbox}[colback=gray!10!white,colframe=gray!75!black,title=\textbf{#1}]
#2
\end{tcolorbox}
}

\newcommand{\correctionbox}[2]{%
\begin{tcolorbox}[colback=red!5!white,colframe=red!75!black,title=\textbf{#1}]
#2
\end{tcolorbox}
}

\newcommand{\newbox}[1]{%
\begin{tcolorbox}[colback=blue!5!white,colframe=blue!75!black,title=\vnewfour]
#1
\end{tcolorbox}
}

%% ============================================================================
%% HYPERREF SETUP
%% ============================================================================
\hypersetup{
    colorlinks=true,
    linkcolor=chapterblue,
    citecolor=thmgreen,
    urlcolor=defblue,
    pdftitle={TKS Formal Mathematical Manual v7.4 - Corrected Edition},
    pdfauthor={TKS Formalization Project}
}

%% ============================================================================
%% DOCUMENT INFO
%% ============================================================================
\title{\Huge\bfseries The TKS Formal Mathematical Manual\\[0.5em]
\Large Version 7.4 --- Corrected Edition\\[1em]
\normalsize Extended Fractals $\cdot$ Domain Semantics $\cdot$ Higher Categories $\cdot$ Noetic Geometry}
\author{TKS Formalization Project}
\date{January 2026}

%% ============================================================================
%% BEGIN DOCUMENT
%% ============================================================================
\begin{document}

\maketitle

\frontmatter

%% ============================================================================
%% CANONICAL AUTHORITY DECLARATION
%% ============================================================================
\chapter*{Canonical Authority Declaration}
\addcontentsline{toc}{chapter}{Canonical Authority Declaration}

\infobox{v7.3 $\to$ v7.4 Continuity Bridge}{
\textbf{Canonical Inheritance Declaration.} This document, TKS v7.4, is a \emph{strict superset} of TKS v7.3, which itself preserves all of v7.2, v7.1, v7.0, and v6.1. All definitions, theorems, axioms, and proofs from previous versions remain valid and unchanged. Version 7.4 introduces four new domain expansions---no modifications to existing canon.
}

\textbf{Canonical Status of Version 7.4.} This document constitutes the canonical and authoritative specification of the TKS Formal Mathematical System. Version 7.4 fully integrates, refines, and supersedes all prior releases in the v3.x--v7.3 development line. Earlier versions remain valuable as historical strata documenting the conceptual evolution of TKS, but no longer serve as normative references for definitions, constructions, or semantics.

The present release establishes the \textbf{definitive ontology of TKS}, consisting of:
\begin{itemize}
    \item The forty canonical Elements (A1--D10)
    \item The ten Noetic operators $\{\nu_0, \ldots, \nu_9\}$
    \item The seven Foundations and their twenty-eight Sub-Foundations
    \item The twenty-two Acquisitions (A0 together with the D, W, and P chains)
\end{itemize}

\textbf{Normative Baseline:} All subsequent work, commentary, or extensions shall be considered extensions of v7.4 and must adopt the definitions and conventions codified herein, unless explicitly superseded by a future canonical release.

\section*{Inheritance Chain}

\begin{description}
    \item[v6.1 Canon:] Complete ontology (4 Worlds, 40 Elements, 10 Noetics, 7 Foundations, 22 Acquisitions), noetic algebra, Tootra arithmetic, Category Noetica, ACBE functor
    
    \item[v7.0 Canon:] RPM Monad, Noetic Fractal Calculus, Dynamic Semantics, Type System, Concurrency, Quantum Algebra foundations
    
    \item[v7.1 Canon:] Natural transformations, Domain theory (dcpos, Scott topology), Extended denotational semantics, Topological noetic space, Multi-goal RPM
    
    \item[v7.2 Canon:] Transfinite fractals (preliminary), Measure theory, Polish spaces, Algorithm W, Compiler architecture, Noetic Hilbert spaces, Spectral theory, Monoidal 2-categories, Topos foundations, Coalgebraic dynamics
    
    \item[v7.3 Canon:] Full Transfinite Noetic Fractals, Extended TKS Language/VM, Complete Quantum Noetics, 2-Category/Topos/Coalgebra Unification
\end{description}

\section*{v7.4 Major Domain Expansions}

\begin{enumerate}
    \item \textbf{MATH Track: Extended Transfinite Fractals}
    \begin{itemize}
        \item Large ordinal hierarchies and inaccessible cardinals
        \item Multi-indexed fractal systems
        \item Fractal spectrum theory
        \item Functorial fractal constructions
    \end{itemize}
    
    \item \textbf{SEMANTICS Track: Domain-Theoretic Semantics}
    \begin{itemize}
        \item Scott domains for noetic computation
        \item Continuous lattices and way-below relation
        \item Fixed-point semantics for recursive noetics
        \item Domain equations and solutions
    \end{itemize}
    
    \item \textbf{HIGHER CATEGORIES Track: $(\infty,1)$-Categories}
    \begin{itemize}
        \item Quasi-categories and Segal spaces
        \item $(\infty,1)$-topoi for noetic sheaves
        \item Noetic homotopy theory
        \item Higher stacks and descent
    \end{itemize}
    
    \item \textbf{GEOMETRY Track: Differential Noetic Geometry}
    \begin{itemize}
        \item Noetic manifolds and tangent structures
        \item Noetic connections and curvature
        \item Geodesics in Idea space
        \item Holonomy and parallel transport
    \end{itemize}
\end{enumerate}

\section*{New Notation in v7.4}

\begin{table}[H]
\centering
\begin{tabular}{lll}
\toprule
\textbf{Symbol} & \textbf{Meaning} & \textbf{Track} \\
\midrule
$\Omega$ & Large ordinal classes & Math \\
$\mathfrak{H}_{\Ord}$ & Ordinal hierarchy & Math \\
$\sigma_\Phi$ & Fractal spectrum & Math \\
$\mathbf{ScottDom}$ & Scott domains & Semantics \\
$\ll$ & Way-below relation & Semantics \\
$\mathsf{lfp}$ & Kleene fixed point & Semantics \\
$(\infty,1)\text{-}\mathbf{Cat}$ & $(\infty,1)$-categories & Higher Categories \\
$(\infty,1)\text{-}\mathbf{Topos}$ & $(\infty,1)$-topoi & Higher Categories \\
$\pi_n^\nu$ & Noetic homotopy groups & Higher Categories \\
$M_\nu$ & Noetic manifold & Geometry \\
$\nabla^\nu$ & Noetic connection & Geometry \\
$R^\nu$ & Noetic curvature & Geometry \\
\bottomrule
\end{tabular}
\end{table}

\tableofcontents

\mainmatter

%% ============================================================================
%% PART I: CANONICAL FOUNDATION
%% ============================================================================
\part{Canonical Foundation (v6.1 Reference)}
\label{part:v61-reference}

\chapter{Ontological Foundations Reference}
\label{ch:foundations-reference}

\infobox{Reference Chapter}{
This chapter provides a condensed reference to the canonical v6.1 foundations. For complete details, see the TKS Formal Mathematical Manual v6.1.
}

\section{Terminology Equivalence for Foundations}
\label{sec:foundation-terminology}

\begin{table}[H]
\centering
\caption{Terminology Equivalence for the Seven Foundations}
\label{tab:foundation-equivalence}
\renewcommand{\arraystretch}{1.2}
\adjustbox{max width=\textwidth}{%
\begin{tabular}{c l l p{5.5cm}}
\toprule
\textbf{Foundation} & \textbf{Metaphysical Name} & \textbf{Canonical v7.4 Label} & \textbf{Formal Principle} \\
\midrule
$\Foundations_1$ & Unity & Association & Coherence and linking of Ideas across Worlds \\
$\Foundations_2$ & Wisdom & Wisdom & Discernment, truth-orientation, structural understanding \\
$\Foundations_3$ & Life & Life & Vitality, persistence, and self-maintenance of patterns \\
$\Foundations_4$ & Companionship & Companionship & Relational bonding, mutual support, shared context \\
$\Foundations_5$ & Power & Power & Capacity to influence, direct, or constrain dynamics \\
$\Foundations_6$ & Material & Wealth & Resources, structure, and embodied stability \\
$\Foundations_7$ & Lust / Continuation & Continuation & Drive to extend, propagate, and iterate patterns \\
\bottomrule
\end{tabular}%
}
\end{table}

\begin{remarkbox*}{Interpretive Notes}{}
The metaphysical labels in the second column reflect the historical and experiential language used in earlier Tootra expositions. The canonical v7.4 labels in the third column capture the corresponding formal principles as they function within the algebraic and categorical TKS calculus. Throughout this manual, the Foundations $\Foundations_1,\dots,\Foundations_7$ are treated as the primary formal objects, with metaphysical names retained for interpretive continuity and cross-referencing to earlier texts.

\textbf{Usage convention}: In formal proofs and definitions, use the canonical v7.4 label. In expository or philosophical discussion, either label is acceptable. The index notation $\Foundations_i$ is universal.
\end{remarkbox*}

\subsection{The 28 Sub-Foundations (Canonical Grid)}
\label{sec:28-subfoundations}

This subsection consolidates the full canonical structure of the \emph{twenty-eight Sub-Foundations} $\Foundations_{mw}$, which refine the seven Foundations $\Foundations_m$ across the four Worlds $w \in \{A,B,C,D\}$.

Each Sub-Foundation $\Foundations_{mw}$ denotes the contextualization of Foundation $\Foundations_m$ within World $w$:
\[
\Foundations_{mw} : \mathrm{Domain} \to \mathrm{Domain}
\]
where the semantic effect of $\Foundations_{mw}$ is to filter or reinterpret a TKS expression with respect to the structural constraints of world-indexed cognition, affect, causality, or embodiment.

\begin{table}[H]
\centering
\caption{Canonical 28 Sub-Foundations}
\label{tab:subfoundations}
\renewcommand{\arraystretch}{1.1}
\adjustbox{max width=\textwidth}{%
\begin{tabular}{c|c|c|p{6.5cm}}
\toprule
$m$ & \textbf{Foundation} & \textbf{World} $w$ & \textbf{Formal Interpretation} \\
\midrule
1 & Unity / Association   & A & Abstract unification of Ideas (A-World coherence). \\
  &                       & B & Cognitive unification: category-level grouping. \\
  &                       & C & Affective synchrony and interpersonal resonance. \\
  &                       & D & Material continuity: object persistence. \\
\midrule
2 & Wisdom / Structure    & A & Archetypal ordering; ideal structural pattern. \\
  &                       & B & Conceptual structuring; mental models. \\
  &                       & C & Emotional evaluation; felt coherence. \\
  &                       & D & Physical constraints; environmental affordances. \\
\midrule
3 & Life / Animation      & A & Essential vitality; causal spark. \\
  &                       & B & Motivation; cognitive impetus to act. \\
  &                       & C & Emotional animation; energizing affect. \\
  &                       & D & Biological / kinetic activation. \\
\midrule
4 & Companionship / Relation & A & Archetypal relational patterns. \\
  &                          & B & Social cognition; theory of mind. \\
  &                          & C & Emotional bonding; empathy dynamics. \\
  &                          & D & Material relations; interaction pathways. \\
\midrule
5 & Power / Expression    & A & Abstract intentionality; divine will-pattern. \\
  &                       & B & Strategic cognition; directed agency. \\
  &                       & C & Emotional assertion; boundary formation. \\
  &                       & D & Physical expression; force deployment. \\
\midrule
6 & Material / Acquisition & A & Ideal pattern of resource flow. \\
  &                        & B & Conceptual economic models. \\
  &                        & C & Emotional valuation; attachment/aversion. \\
  &                        & D & Physical resources; material accumulation. \\
\midrule
7 & Lust / Continuation   & A & Archetypal continuation; pattern perpetuation. \\
  &                       & B & Cognitive projection: planning continuations. \\
  &                       & C & Desire dynamics; emotional momentum. \\
  &                       & D & Biological reproductive / survival drives. \\
\bottomrule
\end{tabular}%
}
\end{table}

\chapter{Core Calculus Overview}
\label{ch:core-calculus}

\section{Core Objects}

\begin{definitionbox*}{The Four Worlds}{worlds}
The TKS ontology is stratified into four \textbf{Worlds}, denoted $\Worlds = \{A, B, C, D\}$:
\begin{itemize}
    \item \textbf{World A} (Archetypal): The realm of pure Ideas and abstract patterns
    \item \textbf{World B} (Mental): The realm of cognition, thought, and mental processes
    \item \textbf{World C} (Emotional): The realm of affect, feeling, and relational dynamics
    \item \textbf{World D} (Physical): The realm of matter, embodiment, and physical action
\end{itemize}
The Worlds are ordered by the ACBE descent: $A \to B \to C \to D$.
\end{definitionbox*}

\begin{definitionbox*}{The Forty Ideas}{ideas}
The set $\Ideas$ consists of 40 \textbf{Ideas}, partitioned by World:
\[
\Ideas = \Ideas_A \sqcup \Ideas_B \sqcup \Ideas_C \sqcup \Ideas_D
\]
where each World-partition contains 10 Ideas:
\begin{itemize}
    \item $\Ideas_A = \{A1, A2, \ldots, A10\}$
    \item $\Ideas_B = \{B1, B2, \ldots, B10\}$
    \item $\Ideas_C = \{C1, C2, \ldots, C10\}$
    \item $\Ideas_D = \{D1, D2, \ldots, D10\}$
\end{itemize}
\end{definitionbox*}

\begin{definitionbox*}{The Ten Noetic Operators}{noetics}
The \textbf{Noetic operators} $\{\nu_0, \nu_1, \ldots, \nu_9\}$ are transformations on Ideas:
\begin{enumerate}[label=(\arabic*)]
    \item $\nu_0$: \textbf{Mind} --- Identity/consciousness operator
    \item $\nu_1$: \textbf{Positive} --- Affirmation/expansion
    \item $\nu_2$: \textbf{Negative} --- Negation/contraction
    \item $\nu_3$: \textbf{Vibration} --- Oscillation/rhythm
    \item $\nu_4$: \textbf{Female} --- Receptive/integrative
    \item $\nu_5$: \textbf{Male} --- Active/differentiating
    \item $\nu_6$: \textbf{Rhythm} --- Cyclic/periodic
    \item $\nu_7$: \textbf{Above/Cause} --- Causal/originating
    \item $\nu_8$: \textbf{Below/Effect} --- Effectual/resulting
    \item $\nu_9$: \textbf{Idea} --- Ideation/conceptualization
\end{enumerate}
\end{definitionbox*}

\section{Algebraic Structure}

\begin{definitionbox*}{Tootra Arithmetic}{tootra-arith}
The \textbf{Tootra arithmetic} operations on Ideas are:
\begin{itemize}
    \item $\Tadd$: Tootra addition (union of noetic content)
    \item $\Tsub$: Tootra subtraction (difference of noetic content)
    \item $\Tmul$: Tootra multiplication (composition of noetic effects)
    \item $\Tdiv$: Tootra division (decomposition of noetic effects)
\end{itemize}
These operations satisfy specific algebraic laws detailed in v6.1.
\end{definitionbox*}

\section{The RPM Monad}

\begin{definitionbox*}{RPM Monad}{rpm-monad}
The \textbf{RPM (Requisite-Prerequisite-Manifestation) Monad} $\RPM$ is defined by:
\begin{itemize}
    \item \textbf{Unit}: $\eta : \Ideas \to \RPM(\Ideas)$ embeds an Idea into the RPM context
    \item \textbf{Multiplication}: $\mu : \RPM(\RPM(\Ideas)) \to \RPM(\Ideas)$ flattens nested RPM computations
    \item \textbf{Bind}: $\rpmbind : \RPM(\Ideas) \times (\Ideas \to \RPM(\Ideas)) \to \RPM(\Ideas)$ sequences RPM computations
\end{itemize}
The RPM Monad encapsulates the prerequisite-checking semantics of TKS goal pursuit.
\end{definitionbox*}

\section{Finite Fractals}

\begin{definitionbox*}{Noetic Fractal}{fractal}
A \textbf{Noetic Fractal} $\Frac = \langle k_1 : k_2 : \cdots : k_n \rangle$ is a finite sequence of noetic operator indices, representing the composition:
\[
\Frac = \nu_{k_n} \circ \nu_{k_{n-1}} \circ \cdots \circ \nu_{k_1}
\]
Fractals act on Ideas: $\Frac(I) = \nu_{k_n}(\nu_{k_{n-1}}(\cdots \nu_{k_1}(I) \cdots))$.
\end{definitionbox*}

\begin{theorembox*}{Fractal Closure}{fractal-closure}
The set of all finite Noetic Fractals forms a monoid under concatenation, with the empty fractal $\langle \rangle$ (identity) as the unit.
\end{theorembox*}

%% ============================================================================
%% PART II: HIGHER CATEGORIES TRACK (v7.4 New)
%% ============================================================================
\part{Higher Categories Track}
\label{part:higher-categories}

%%%%%%%%%%%%%%%%%%%%%%%%%%%%%%%%%%%%%%%%%%%%%%%%%%%%%%%%%%%%%%%%%%%%%%%%%%%%%%%
%% TKS v7.4 HIGHER CATEGORIES TRACK
%% $(\infty,1)$-Categories, Quasi-Categories, and Noetic Homotopy Theory
%% Part of TKS v7.4 - December 2025
%% ADDITIVE to v7.3 canon
%%%%%%%%%%%%%%%%%%%%%%%%%%%%%%%%%%%%%%%%%%%%%%%%%%%%%%%%%%%%%%%%%%%%%%%%%%%%%%%

%% ============================================================================
%% CHAPTER 1: QUASI-CATEGORIES AND NOETIC SPACES
%% ============================================================================

\chapter{Quasi-Categories and Noetic Spaces}
\label{ch:quasi-categories}

\begin{tcolorbox}[colback=blue!5!white,colframe=blue!75!black,title=\vnewfour]
This chapter introduces quasi-categories as a model for $(\infty,1)$-categories
and constructs the noetic quasi-category $\mathcal{N}_\infty$.
\end{tcolorbox}

\section{Simplicial Preliminaries}

\begin{definition}[Simplex Category]
\label{def:simplex-cat}
The \textbf{simplex category} $\Delta$ has:
\begin{itemize}
  \item Objects: Finite ordinals $[n] = \{0, 1, \ldots, n\}$ for $n \geq 0$
  \item Morphisms: Order-preserving maps
\end{itemize}
\end{definition}

\begin{definition}[Simplicial Set]
\label{def:simplicial-set}
A \textbf{simplicial set} is a functor $X : \Delta^{\mathrm{op}} \to \mathbf{Set}$.
We write $X_n := X([n])$ for the set of $n$-simplices.
\end{definition}

\begin{definition}[Horn]
\label{def:horn}
For $0 \leq k \leq n$, the \textbf{$k$-horn} $\Lambda^n_k \subset \Delta^n$ is the
simplicial subset generated by all faces except the $k$-th face.
\end{definition}

\section{Definition of Quasi-Categories}

\begin{definition}[Quasi-Category]
\label{def:quasi-category}
A simplicial set $\mathcal{C}$ is a \textbf{quasi-category} (or $\infty$-category)
if it has the \emph{inner horn filling property}: for all $0 < k < n$, every map
$\Lambda^n_k \to \mathcal{C}$ extends to $\Delta^n \to \mathcal{C}$.
\[
\begin{tikzcd}
\Lambda^n_k \arrow[r] \arrow[d, hook] & \mathcal{C} \\
\Delta^n \arrow[ur, dashed] &
\end{tikzcd}
\]
\end{definition}

\begin{remark}[Interpretation]
In a quasi-category:
\begin{itemize}
  \item 0-simplices are objects
  \item 1-simplices are morphisms
  \item 2-simplices witness composition (up to homotopy)
  \item Higher simplices encode coherence data
\end{itemize}
\end{remark}

\section{The Noetic Quasi-Category $\mathcal{N}_\infty$}

\begin{construction}[Noetic Quasi-Category]
\label{constr:noetic-quasi-cat}
The \textbf{noetic quasi-category} $\mathcal{N}_\infty$ is defined by:
\begin{itemize}
  \item $(\mathcal{N}_\infty)_0 = \Ideas$ (Ideas are objects)
  \item $(\mathcal{N}_\infty)_1 = \{(\iota, \iota', \Phi) : \Phi \text{ is a noetic path from } \iota \text{ to } \iota'\}$
  \item $(\mathcal{N}_\infty)_n$ consists of $n$-chains of noetic paths with coherent homotopies
\end{itemize}
\end{construction}

\begin{theorem}[Noetic Quasi-Category is Well-Defined]
\label{thm:noetic-qcat-well-defined}
$\mathcal{N}_\infty$ is a quasi-category. Inner horn fillings correspond to
fractal interpolation between noetic paths.
\end{theorem}

\section{Mapping Spaces and Hom-Objects}

\begin{definition}[Mapping Space]
\label{def:mapping-space}
For $\iota, \iota' \in \mathcal{N}_\infty$, the \textbf{mapping space} is:
\[
\mathrm{Map}_{\mathcal{N}_\infty}(\iota, \iota') := \{\sigma \in (\mathcal{N}_\infty)_\bullet : d_0\sigma = \iota', d_n\sigma = \iota\}
\]
This is a Kan complex.
\end{definition}

\begin{proposition}[Noetic Mapping Spaces]
\label{prop:noetic-mapping}
The mapping space $\mathrm{Map}_{\mathcal{N}_\infty}(\iota, \iota')$ has:
\begin{enumerate}[label=(\roman*)]
  \item 0-simplices: noetic morphisms $\iota \to \iota'$
  \item 1-simplices: homotopies between morphisms
  \item $n$-simplices: higher homotopy data
\end{enumerate}
\end{proposition}

\section{Kan Complexes and Fibrancy}

\begin{definition}[Kan Complex]
\label{def:kan-complex}
A simplicial set $K$ is a \textbf{Kan complex} if all horns (inner and outer)
can be filled.
\end{definition}

\begin{proposition}[Mapping Spaces are Kan]
\label{prop:mapping-kan}
For any quasi-category $\mathcal{C}$ and objects $x, y$, the mapping space
$\mathrm{Map}_\mathcal{C}(x, y)$ is a Kan complex.
\end{proposition}

%% ============================================================================
%% CHAPTER 2: SEGAL SPACES FOR NOETICS
%% ============================================================================

\chapter{Segal Spaces for Noetics}
\label{ch:segal-spaces}

\begin{tcolorbox}[colback=blue!5!white,colframe=blue!75!black,title=\vnewfour]
This chapter develops Segal spaces as an alternative model equivalent to
quasi-categories.
\end{tcolorbox}

\section{Segal Spaces}

\begin{definition}[Segal Space]
\label{def:segal-space}
A \textbf{Segal space} is a simplicial space $X : \Delta^{\mathrm{op}} \to \mathbf{Space}$
satisfying the \emph{Segal condition}: the map
\[
X_n \to X_1 \times_{X_0} X_1 \times_{X_0} \cdots \times_{X_0} X_1
\]
is a weak equivalence for all $n \geq 2$.
\end{definition}

\begin{definition}[Complete Segal Space]
\label{def:complete-segal}
A Segal space $X$ is \textbf{complete} (or \textbf{Rezk complete}) if the map
\[
X_0 \to X_1^{\mathrm{eq}}
\]
is a weak equivalence, where $X_1^{\mathrm{eq}}$ is the space of equivalences.
\end{definition}

\section{Noetic Segal Space}

\begin{construction}[Noetic Segal Space]
\label{constr:noetic-segal}
The \textbf{noetic Segal space} $\mathcal{N}^{\mathrm{Seg}}$ has:
\begin{itemize}
  \item $\mathcal{N}^{\mathrm{Seg}}_0 = |\Ideas|$ (geometric realization of Idea space)
  \item $\mathcal{N}^{\mathrm{Seg}}_1 = $ space of noetic paths
  \item $\mathcal{N}^{\mathrm{Seg}}_n = $ space of composable $n$-tuples
\end{itemize}
\end{construction}

\begin{theorem}[Noetic Segal Condition]
\label{thm:noetic-segal-condition}
The noetic Segal space satisfies the Segal condition via fractal composition.
\end{theorem}

\section{Equivalence with Quasi-Categories}

\begin{theorem}[Segal-Quasi-Category Equivalence]
\label{thm:segal-qcat-equiv}
There is a Quillen equivalence between:
\begin{itemize}
  \item Complete Segal spaces
  \item Quasi-categories
\end{itemize}
Under this equivalence, $\mathcal{N}^{\mathrm{Seg}} \simeq \mathcal{N}_\infty$.
\end{theorem}

%% ============================================================================
%% CHAPTER 3: $(\infty,1)$-CATEGORIES OF IDEAS
%% ============================================================================

\chapter{$(\infty,1)$-Categories of Ideas}
\label{ch:infty-ideas}

\begin{tcolorbox}[colback=blue!5!white,colframe=blue!75!black,title=\vnewfour]
This chapter develops the full $(\infty,1)$-category $\mathbf{Idea}_\infty$.
\end{tcolorbox}

\section{The $(\infty,1)$-Category $\mathbf{Idea}_\infty$}

\begin{definition}[$(\infty,1)$-Category of Ideas]
\label{def:idea-infty}
The \textbf{$(\infty,1)$-category of Ideas} $\mathbf{Idea}_\infty$ extends the
v6.1 category $\Noetica$ with:
\begin{itemize}
  \item Objects: Ideas $\iota \in \Ideas$
  \item 1-morphisms: Noetic transformations
  \item 2-morphisms: Homotopies between transformations
  \item $n$-morphisms: Higher homotopy coherences
\end{itemize}
\end{definition}

\begin{proposition}[World Filtration]
\label{prop:world-filtration}
$\mathbf{Idea}_\infty$ admits a filtration by Worlds:
\[
\mathbf{Idea}_\infty^{(A)} \hookrightarrow \mathbf{Idea}_\infty^{(B)} \hookrightarrow
\mathbf{Idea}_\infty^{(C)} \hookrightarrow \mathbf{Idea}_\infty^{(D)} = \mathbf{Idea}_\infty
\]
\end{proposition}

\section{Higher Morphisms Between Ideas}

\begin{definition}[Higher Noetic Morphism]
\label{def:higher-morphism}
An \textbf{$n$-morphism} in $\mathbf{Idea}_\infty$ is a map
\[
\alpha : \Delta^n \to \mathbf{Idea}_\infty
\]
with boundary conditions specifying source and target $(n-1)$-morphisms.
\end{definition}

\begin{proposition}[Truncation Recovers $\Noetica$]
\label{prop:truncation}
The 1-truncation $\tau_{\leq 1}\mathbf{Idea}_\infty$ is equivalent to the ordinary
category $\Noetica$ from v6.1.
\end{proposition}

\section{Adjunctions in the $\infty$-Setting}

\begin{definition}[$\infty$-Adjunction]
\label{def:infty-adjunction}
An \textbf{$\infty$-adjunction} between $(\infty,1)$-categories $\mathcal{C}$ and $\mathcal{D}$
consists of functors $F : \mathcal{C} \rightleftarrows \mathcal{D} : G$ with a natural
equivalence of mapping spaces:
\[
\mathrm{Map}_\mathcal{D}(F(x), y) \simeq \mathrm{Map}_\mathcal{C}(x, G(y))
\]
\end{definition}

\begin{theorem}[Noetic Adjunctions]
\label{thm:noetic-adjunctions}
The following are $\infty$-adjunctions on $\mathbf{Idea}_\infty$:
\begin{enumerate}[label=(\roman*)]
  \item \textbf{ACBE adjunction}: World descent $\dashv$ World ascent
  \item \textbf{Fractal adjunction}: Fractal application $\dashv$ Fractal extraction
  \item \textbf{RPM adjunction}: Free RPM monad $\dashv$ Forgetful
\end{enumerate}
\end{theorem}

%% ============================================================================
%% CHAPTER 4: $(\infty,1)$-TOPOI AND NOETIC SHEAVES
%% ============================================================================

\chapter{$(\infty,1)$-Topoi and Noetic Sheaves}
\label{ch:infty-topoi}

\begin{tcolorbox}[colback=blue!5!white,colframe=blue!75!black,title=\vnewfour]
This chapter develops the noetic $\infty$-topos $\mathbf{Noe}_\infty$.
\end{tcolorbox}

\section{$(\infty,1)$-Topoi}

\begin{definition}[Presentable $(\infty,1)$-Category]
\label{def:presentable}
An $(\infty,1)$-category is \textbf{presentable} if it is accessible and admits
all small colimits.
\end{definition}

\begin{definition}[$(\infty,1)$-Topos]
\label{def:infty-topos}
An \textbf{$(\infty,1)$-topos} is a presentable $(\infty,1)$-category $\mathcal{E}$
that is left exact localization of a presheaf category:
\[
\mathcal{E} \simeq L\mathbf{PSh}(\mathcal{C})
\]
for some small $(\infty,1)$-category $\mathcal{C}$.
\end{definition}

\section{The Noetic $\infty$-Topos}

\begin{definition}[Noetic $\infty$-Topos]
\label{def:noetic-infty-topos}
The \textbf{noetic $\infty$-topos} $\InftyTopos_\nu$ is defined as:
\[
\InftyTopos_\nu := \Sh_\infty(\Ideas, \tau_{\mathrm{noetic}})
\]
the $(\infty,1)$-category of $\infty$-sheaves on the noetic site.
\end{definition}

\begin{theorem}[$\InftyTopos_\nu$ is an $\infty$-Topos]
\label{thm:noe-infty-topos}
$\InftyTopos_\nu$ satisfies the axioms of an $(\infty,1)$-topos.
\end{theorem}

\section{Descent and Higher Sheaf Conditions}

\begin{definition}[$\infty$-Descent]
\label{def:infty-descent}
A presheaf $F : \mathcal{C}^{\mathrm{op}} \to \mathbf{Space}$ satisfies \textbf{$\infty$-descent}
for a covering sieve $S$ if the natural map
\[
F(U) \to \lim_{\Delta} F(S_\bullet)
\]
is an equivalence, where $S_\bullet$ is the \v{C}ech nerve.
\end{definition}

\begin{theorem}[Noetic Descent]
\label{thm:noetic-descent}
The noetic topology satisfies effective $\infty$-descent: an $\infty$-presheaf
is a sheaf iff it satisfies descent for all noetic coverings.
\end{theorem}

\section{$\infty$-Giraud Axioms}

\begin{theorem}[$\infty$-Giraud Axioms]
\label{thm:infty-giraud}
An $(\infty,1)$-category $\mathcal{E}$ is an $\infty$-topos iff:
\begin{enumerate}[label=(G\arabic*)]
  \item $\mathcal{E}$ has all small colimits
  \item Colimits are universal
  \item Coproducts are disjoint
  \item Every groupoid object is effective
  \item There exists a set of generators
\end{enumerate}
\end{theorem}

\begin{corollary}[$\InftyTopos_\nu$ Satisfies Giraud Axioms]
\label{cor:noe-giraud}
$\InftyTopos_\nu$ satisfies all $\infty$-Giraud axioms.
\end{corollary}

%% ============================================================================
%% CHAPTER 5: NOETIC HOMOTOPY GROUPS
%% ============================================================================

\chapter{Noetic Homotopy Groups}
\label{ch:noetic-homotopy}

\begin{tcolorbox}[colback=blue!5!white,colframe=blue!75!black,title=\vnewfour]
This chapter develops noetic homotopy groups $\NoeticHomotopy$ and related
invariants.
\end{tcolorbox}

\section{Definition of Noetic Homotopy Groups}

\begin{definition}[Noetic Homotopy Groups]
\label{def:noetic-homotopy-groups}
For an Idea $\iota \in \Ideas$ and $n \geq 0$, the \textbf{$n$-th noetic homotopy group} is:
\[
\NoeticHomotopy(\iota) := \pi_n(\mathrm{Map}_{\mathcal{N}_\infty}(\iota, \iota), \mathrm{id}_\iota)
\]
\end{definition}

\begin{proposition}[Properties]
\label{prop:homotopy-properties}
\begin{enumerate}[label=(\roman*)]
  \item $\pi^\nu_0(\iota)$ classifies connected components of the Idea's automorphism space
  \item $\pi^\nu_1(\iota)$ is the fundamental group of automorphisms
  \item $\pi^\nu_n(\iota)$ for $n \geq 2$ are abelian
\end{enumerate}
\end{proposition}

\section{Noetic Whitehead Theorem}

\begin{theorem}[Noetic Whitehead Theorem]
\label{thm:noetic-whitehead}
A morphism $f : \iota \to \iota'$ in $\mathcal{N}_\infty$ is an equivalence iff
it induces isomorphisms on all noetic homotopy groups:
\[
f_* : \NoeticHomotopy(\iota) \xrightarrow{\cong} \NoeticHomotopy(\iota')
\]
for all $n \geq 0$.
\end{theorem}

\section{Postnikov Towers of Ideas}

\begin{definition}[Postnikov Tower]
\label{def:postnikov-tower}
The \textbf{Postnikov tower} of an Idea $\iota$ is a sequence:
\[
\cdots \to P_n(\iota) \to P_{n-1}(\iota) \to \cdots \to P_0(\iota)
\]
where $P_n(\iota)$ is the $n$-truncation with $\pi^\nu_k(P_n(\iota)) = 0$ for $k > n$.
\end{definition}

\begin{theorem}[Postnikov Convergence]
\label{thm:postnikov-conv}
\[
\iota \simeq \lim_n P_n(\iota)
\]
The Idea is recovered as the limit of its Postnikov tower.
\end{theorem}

\section{Homotopy Limits and Colimits}

\begin{definition}[Homotopy Limit]
\label{def:homotopy-limit}
The \textbf{homotopy limit} of a diagram $D : J \to \mathcal{N}_\infty$ is:
\[
\mathrm{holim}_J D := \mathrm{Map}(\mathcal{N}(J), \mathcal{N}_\infty) \times_{\mathrm{Map}(\partial\mathcal{N}(J), \mathcal{N}_\infty)} \{D\}
\]
\end{definition}

\begin{theorem}[Milnor Exact Sequence]
\label{thm:milnor}
For a tower of Ideas $\cdots \to \iota_2 \to \iota_1 \to \iota_0$:
\[
0 \to \lim^1 \pi^\nu_{n+1}(\iota_k) \to \pi^\nu_n(\mathrm{holim}_k \iota_k) \to \lim \pi^\nu_n(\iota_k) \to 0
\]
\end{theorem}

%% ============================================================================
%% CHAPTER 6: HIGHER STACKS AND MODULI
%% ============================================================================

\chapter{Higher Stacks and Moduli}
\label{ch:higher-stacks}

\begin{tcolorbox}[colback=blue!5!white,colframe=blue!75!black,title=\vnewfour]
This chapter develops noetic higher stacks and moduli problems.
\end{tcolorbox}

\section{Higher Stacks}

\begin{definition}[Higher Stack]
\label{def:higher-stack}
A \textbf{higher stack} (or $\infty$-stack) on the noetic site is a functor
\[
\mathcal{F} : \Ideas^{\mathrm{op}} \to \mathbf{Space}
\]
satisfying $\infty$-descent for all noetic coverings.
\end{definition}

\begin{definition}[Category of Noetic Stacks]
\label{def:stack-category}
The \textbf{category of noetic stacks} is:
\[
\HigherStack := \Sh_\infty(\Ideas, \tau_{\mathrm{noetic}})
\]
\end{definition}

\section{Moduli Problems for Ideas}

\begin{definition}[Moduli Stack of Ideas]
\label{def:moduli-ideas}
The \textbf{moduli stack of Ideas} $\mathcal{M}_{\Ideas}$ assigns to each Idea $\iota$
the space of Ideas ``over'' $\iota$:
\[
\mathcal{M}_{\Ideas}(\iota) := \mathrm{Map}_{\mathcal{N}_\infty}(\iota, \mathbf{Idea}_\infty)
\]
\end{definition}

\begin{definition}[Moduli of Fractals]
\label{def:moduli-fractals}
The \textbf{moduli stack of fractals} $\mathcal{M}_\Phi$ parametrizes fractal
structures on Ideas:
\[
\mathcal{M}_\Phi(\iota) := \{\Phi : \Phi \text{ is a fractal structure on } \iota\}
\]
\end{definition}

\section{Descent Conditions}

\begin{theorem}[ACBE Descent]
\label{thm:acbe-descent}
ACBE morphisms satisfy effective descent: for any ACBE covering
$\{U_\alpha \to X\}$, the natural functor
\[
\HigherStack(X) \to \mathrm{holim}_\Delta \HigherStack(U_\bullet)
\]
is an equivalence.
\end{theorem}

\section{Classifying Stacks}

\begin{definition}[Classifying Stack]
\label{def:classifying-stack}
The \textbf{classifying stack} $B\mathrm{Aut}(\iota)$ of an Idea $\iota$ represents
the functor:
\[
\iota' \mapsto \{\text{$\iota$-torsors over } \iota'\}
\]
\end{definition}

\begin{theorem}[Classification]
\label{thm:classification}
For connected $\iota'$:
\[
\mathrm{Map}(\iota', B\mathrm{Aut}(\iota)) \simeq \mathrm{Tors}_\iota(\iota')
\]
The mapping space classifies $\iota$-torsors.
\end{theorem}

%% ============================================================================
%% CHAPTER 7: SYNTHESIS AND CONNECTIONS
%% ============================================================================

\chapter{Synthesis and Connections}
\label{ch:synthesis}

\begin{tcolorbox}[colback=green!5!white,colframe=green!75!black,title=\textbf{Integration}]
This chapter connects the higher categorical framework to other v7.4 tracks.
\end{tcolorbox}

\section{Connection to Transfinite Fractals (Part V)}

\begin{proposition}[Fractal Functor to $\infty$-Categories]
\label{prop:fractal-infty}
The fractal functor $\FractalFunctor : \Ord \to \mathbf{Cat}$ extends to an
$\infty$-functor:
\[
\FractalFunctor_\infty : \Ord \to \InftyCat
\]
\end{proposition}

\section{Connection to Domain Semantics (Part VI)}

\begin{proposition}[Domain-$\infty$-Category Correspondence]
\label{prop:domain-infty}
Scott domains embed into $\infty$-groupoids via the nerve construction:
\[
\mathcal{N} : \ScottDom \hookrightarrow \InftyCat
\]
\end{proposition}

\section{Connection to v7.3 Meta Track}

\begin{theorem}[2-Category Upgrade]
\label{thm:2cat-upgrade}
The v7.3 2-categorical structures embed into $(\infty,1)$-categories:
\[
\TwoCat \hookrightarrow \InftyCat
\]
preserving adjunctions and limits.
\end{theorem}

\section{Summary Theorem}

\begin{theorem}[TKS Higher Categorical Foundations]
\label{thm:summary}
The noetic $(\infty,1)$-category $\mathbf{Idea}_\infty$ satisfies:
\begin{enumerate}[label=(\roman*)]
  \item It is an $\infty$-topos
  \item It admits all homotopy limits and colimits
  \item Noetic homotopy groups provide complete invariants
  \item Higher stacks classify moduli of noetic structures
  \item All v7.0--v7.3 structures embed faithfully
\end{enumerate}
\end{theorem}

\begin{tcolorbox}[colback=green!5!white,colframe=green!75!black,title=\textbf{Cross-Track Connections}]
\textbf{To Part V (Math)}: Ordinal-indexed $\infty$-categories via $\FractalFunctor_\infty$.

\textbf{To Part VI (Semantics)}: Domain-theoretic models as $\infty$-groupoids.

\textbf{To Part VIII (Geometry)}: Noetic manifolds as objects in a geometric
$\infty$-topos.
\end{tcolorbox}


%% ============================================================================
%% PART III: DIFFERENTIAL NOETIC GEOMETRY TRACK (v7.4 New)
%% ============================================================================
\part{Differential Noetic Geometry Track}
\label{part:geometry}

%%%%%%%%%%%%%%%%%%%%%%%%%%%%%%%%%%%%%%%%%%%%%%%%%%%%%%%%%%%%%%%%%%%%%%%%%%%%%%%
%% TKS v7.4 DIFFERENTIAL NOETIC GEOMETRY TRACK
%% Noetic Manifolds, Connections, Curvature, Geodesics, and Holonomy
%% Part of TKS v7.4 - December 2025
%% ADDITIVE to v7.3 canon
%%%%%%%%%%%%%%%%%%%%%%%%%%%%%%%%%%%%%%%%%%%%%%%%%%%%%%%%%%%%%%%%%%%%%%%%%%%%%%%

%% ============================================================================
%% CHAPTER 1: NOETIC MANIFOLDS
%% ============================================================================

\chapter{Noetic Manifolds}
\label{ch:noetic-manifolds}

\begin{tcolorbox}[colback=blue!5!white,colframe=blue!75!black,title=\vnewfour]
This chapter introduces the concept of noetic manifolds as smooth spaces whose
points represent idea-states and whose differential structure encodes the local
geometry of ideation.
\end{tcolorbox}

\section{Definition of Noetic Manifolds}

\begin{definition}[Noetic Manifold]
\label{def:noetic-manifold}
A \textbf{noetic manifold} $\NoeticManifold$ is a tuple $(M, \mathcal{A}, \nu)$ where:
\begin{enumerate}[label=(\roman*)]
    \item $(M, \mathcal{A})$ is a smooth manifold with atlas $\mathcal{A}$,
    \item $\nu : M \to \Ideas$ is a smooth map assigning to each point $p \in M$
          an idea $\nu(p) \in \Ideas$,
    \item The fibers of $\nu$ carry a noetic structure compatible with the
          smooth structure.
\end{enumerate}
\end{definition}

\begin{definition}[Noetic Dimension]
\label{def:noetic-dimension}
The \textbf{noetic dimension} of $\NoeticManifold$ at $p \in M$ is:
\[
    \dim_\nu(p) := \dim(T_pM) + \dim(\ker(d\nu_p)^\perp)
\]
where $d\nu_p : T_pM \to T_{\nu(p)}\Ideas$ is the differential of $\nu$.
\end{definition}

\begin{example}[Idea Space as Noetic Manifold]
\label{ex:idea-space-manifold}
The idea space $\Ideas$ itself is a noetic manifold with $\nu = \mathrm{id}_\Ideas$.
Local coordinates $(\iota^1, \ldots, \iota^n)$ represent components of an idea.
\end{example}

\section{Tangent and Cotangent Bundles}

\begin{definition}[Noetic Tangent Bundle]
\label{def:tangent-bundle}
The \textbf{noetic tangent bundle} $\TangentBundle$ is:
\[
    \TangentBundle := \bigsqcup_{p \in M} T_pM
\]
with projection $\pi : \TangentBundle \to M$ and smooth structure induced from $M$.
A \textbf{noetic tangent vector} $v \in T_pM$ represents an infinitesimal
direction of ideation at $p$.
\end{definition}

\begin{definition}[Noetic Cotangent Bundle]
\label{def:cotangent-bundle}
The \textbf{noetic cotangent bundle} $\CotangentBundle$ is:
\[
    \CotangentBundle := \bigsqcup_{p \in M} T_p^*M
\]
Elements $\omega \in T_p^*M$ are \textbf{noetic 1-forms} representing linear
functionals on infinitesimal ideations.
\end{definition}

\begin{proposition}[Noetic Phase Space]
\label{prop:phase-space}
The cotangent bundle $\CotangentBundle$ carries a canonical symplectic structure
$\omega_\nu = d\theta_\nu$ where $\theta_\nu$ is the tautological 1-form.
This makes $\CotangentBundle$ a \textbf{noetic phase space}.
\end{proposition}

\section{Noetic Vector Fields}

\begin{definition}[Noetic Vector Field]
\label{def:noetic-vector-field}
A \textbf{noetic vector field} $X$ on $\NoeticManifold$ is a smooth section
$X : M \to \TangentBundle$ such that $\pi \circ X = \mathrm{id}_M$.
In local coordinates:
\[
    X = X^i(x) \frac{\partial}{\partial x^i}
\]
\end{definition}

\begin{definition}[Noetic Lie Bracket]
\label{def:lie-bracket}
For noetic vector fields $X, Y$, the \textbf{Lie bracket} is:
\[
    [X, Y] = XY - YX = \left(X^j \frac{\partial Y^i}{\partial x^j}
    - Y^j \frac{\partial X^i}{\partial x^j}\right) \frac{\partial}{\partial x^i}
\]
This encodes the non-commutativity of infinitesimal ideation flows.
\end{definition}

\begin{theorem}[Noetic Lie Algebra]
\label{thm:noetic-lie-algebra}
The space $\mathfrak{X}(\NoeticManifold)$ of noetic vector fields forms an
infinite-dimensional Lie algebra over $\mathbb{R}$ under the Lie bracket.
\end{theorem}

%% ============================================================================
%% CHAPTER 2: THE NOETIC METRIC
%% ============================================================================

\chapter{The Noetic Metric}
\label{ch:noetic-metric}

\begin{tcolorbox}[colback=blue!5!white,colframe=blue!75!black,title=\vnewfour]
This chapter develops the noetic metric tensor, which measures distances and
angles in idea space, enabling quantitative analysis of conceptual proximity.
\end{tcolorbox}

\section{Definition of the Noetic Metric}

\begin{definition}[Noetic Metric]
\label{def:noetic-metric}
A \textbf{noetic metric} on $\NoeticManifold$ is a smooth assignment
$p \mapsto \NoeticMetric_p$ of a positive-definite symmetric bilinear form
$\NoeticMetric_p : T_pM \times T_pM \to \mathbb{R}$ for each $p \in M$.
In local coordinates:
\[
    \NoeticMetric = g^\nu_{ij}(x) \, dx^i \otimes dx^j
\]
\end{definition}

\begin{definition}[Noetic Inner Product]
\label{def:noetic-inner-product}
For tangent vectors $u, v \in T_pM$, the \textbf{noetic inner product} is:
\[
    \langle u, v \rangle_\nu := \NoeticMetric_p(u, v) = g^\nu_{ij} u^i v^j
\]
\end{definition}

\begin{definition}[Noetic Norm]
\label{def:noetic-norm}
The \textbf{noetic norm} of a tangent vector $v \in T_pM$ is:
\[
    \|v\|_\nu := \sqrt{\langle v, v \rangle_\nu} = \sqrt{g^\nu_{ij} v^i v^j}
\]
\end{definition}

\section{Noetic Distances and Lengths}

\begin{definition}[Length of Noetic Curve]
\label{def:curve-length}
For a smooth curve $\gamma : [a, b] \to \NoeticManifold$, the \textbf{noetic length} is:
\[
    L_\nu(\gamma) := \int_a^b \|\dot{\gamma}(t)\|_\nu \, dt
    = \int_a^b \sqrt{g^\nu_{ij}(\gamma(t)) \dot{\gamma}^i(t) \dot{\gamma}^j(t)} \, dt
\]
\end{definition}

\begin{definition}[Noetic Distance]
\label{def:noetic-distance}
The \textbf{noetic distance} between points $p, q \in M$ is:
\[
    d_\nu(p, q) := \inf_{\gamma} L_\nu(\gamma)
\]
where the infimum is over all smooth curves $\gamma$ from $p$ to $q$.
\end{definition}

\begin{theorem}[Noetic Metric Space]
\label{thm:noetic-metric-space}
$(\NoeticManifold, d_\nu)$ is a metric space. The metric topology coincides
with the manifold topology.
\end{theorem}

\section{Information-Theoretic Noetic Metrics}

\begin{definition}[Fisher Information Metric]
\label{def:fisher-metric}
When $\NoeticManifold$ parameterizes probability distributions over ideas,
the \textbf{Fisher information metric} is:
\[
    g^\nu_{ij}(\theta) := \mathbb{E}_{p_\theta}\left[
    \frac{\partial \log p_\theta}{\partial \theta^i}
    \frac{\partial \log p_\theta}{\partial \theta^j}
    \right]
\]
where $p_\theta$ is the probability distribution parameterized by $\theta$.
\end{definition}

\begin{proposition}[Information Monotonicity]
\label{prop:info-monotonicity}
The Fisher information metric satisfies the \textbf{information monotonicity}
property: for any stochastic map $T$, the induced metric on the image is
bounded by the original metric.
\end{proposition}

\begin{example}[Kullback-Leibler Distance]
\label{ex:kl-distance}
The geodesic distance under the Fisher metric is related to the
Kullback-Leibler divergence:
\[
    d_\nu(\theta, \theta')^2 \approx 2 D_{\mathrm{KL}}(p_\theta \| p_{\theta'})
\]
for nearby points $\theta, \theta'$.
\end{example}

%% ============================================================================
%% CHAPTER 3: NOETIC CONNECTIONS
%% ============================================================================

\chapter{Noetic Connections}
\label{ch:noetic-connections}

\begin{tcolorbox}[colback=blue!5!white,colframe=blue!75!black,title=\vnewfour]
This chapter introduces noetic connections, which provide a way to compare
tangent vectors at different points and define parallel transport of ideas.
\end{tcolorbox}

\section{Definition of Noetic Connection}

\begin{definition}[Noetic Connection]
\label{def:noetic-connection}
A \textbf{noetic connection} $\NoeticConnection$ on $\NoeticManifold$ is a map
\[
    \NoeticConnection : \mathfrak{X}(M) \times \mathfrak{X}(M) \to \mathfrak{X}(M),
    \quad (X, Y) \mapsto \NoeticConnection_X Y
\]
satisfying:
\begin{enumerate}[label=(\alph*)]
    \item $\mathbb{R}$-linearity in $Y$: $\NoeticConnection_X(aY + bZ) = a\NoeticConnection_X Y + b\NoeticConnection_X Z$
    \item $C^\infty(M)$-linearity in $X$: $\NoeticConnection_{fX} Y = f\NoeticConnection_X Y$
    \item Leibniz rule: $\NoeticConnection_X(fY) = X(f)Y + f\NoeticConnection_X Y$
\end{enumerate}
\end{definition}

\begin{definition}[Christoffel Symbols]
\label{def:christoffel}
In local coordinates, the connection is determined by \textbf{Christoffel symbols}:
\[
    \NoeticConnection_{\partial_i} \partial_j = \Gamma^\nu{}^k_{ij} \partial_k
\]
The full covariant derivative is:
\[
    \NoeticConnection_X Y = X^i \left(\frac{\partial Y^k}{\partial x^i}
    + \Gamma^\nu{}^k_{ij} Y^j\right) \partial_k
\]
\end{definition}

\section{The Levi-Civita Connection}

\begin{definition}[Metric Compatibility]
\label{def:metric-compatible}
A connection $\NoeticConnection$ is \textbf{metric compatible} with $\NoeticMetric$ if:
\[
    X\langle Y, Z \rangle_\nu = \langle \NoeticConnection_X Y, Z \rangle_\nu
    + \langle Y, \NoeticConnection_X Z \rangle_\nu
\]
for all vector fields $X, Y, Z$.
\end{definition}

\begin{definition}[Torsion-Free]
\label{def:torsion-free}
A connection is \textbf{torsion-free} if:
\[
    \NoeticConnection_X Y - \NoeticConnection_Y X = [X, Y]
\]
Equivalently, $\Gamma^\nu{}^k_{ij} = \Gamma^\nu{}^k_{ji}$ (symmetric in lower indices).
\end{definition}

\begin{theorem}[Fundamental Theorem of Noetic Geometry]
\label{thm:levi-civita}
For any noetic metric $\NoeticMetric$, there exists a unique connection
$\NoeticConnection$ that is both metric compatible and torsion-free.
This is the \textbf{Levi-Civita connection}.
\end{theorem}

\begin{proposition}[Christoffel Symbol Formula]
\label{prop:christoffel-formula}
The Christoffel symbols of the Levi-Civita connection are:
\[
    \Gamma^\nu{}^k_{ij} = \frac{1}{2} g^{\nu kl}\left(
    \frac{\partial g^\nu_{il}}{\partial x^j}
    + \frac{\partial g^\nu_{jl}}{\partial x^i}
    - \frac{\partial g^\nu_{ij}}{\partial x^l}
    \right)
\]
where $g^{\nu kl}$ is the inverse metric.
\end{proposition}

\section{Covariant Derivatives of Tensor Fields}

\begin{definition}[Covariant Derivative of 1-Forms]
\label{def:cov-deriv-forms}
For a 1-form $\omega$, the covariant derivative is:
\[
    (\NoeticConnection_X \omega)(Y) = X(\omega(Y)) - \omega(\NoeticConnection_X Y)
\]
In components:
\[
    \NoeticConnection_i \omega_j = \frac{\partial \omega_j}{\partial x^i}
    - \Gamma^\nu{}^k_{ij} \omega_k
\]
\end{definition}

\begin{definition}[Covariant Derivative of Tensors]
\label{def:cov-deriv-tensors}
For a $(r, s)$-tensor $T$, the covariant derivative adds a Christoffel
term for each index:
\[
    \NoeticConnection_k T^{i_1 \cdots i_r}{}_{j_1 \cdots j_s} =
    \frac{\partial T^{i_1 \cdots i_r}{}_{j_1 \cdots j_s}}{\partial x^k}
    + \sum_a \Gamma^\nu{}^{i_a}_{km} T^{i_1 \cdots m \cdots i_r}{}_{j_1 \cdots j_s}
    - \sum_b \Gamma^\nu{}^m_{kj_b} T^{i_1 \cdots i_r}{}_{j_1 \cdots m \cdots j_s}
\]
\end{definition}

%% ============================================================================
%% CHAPTER 4: NOETIC CURVATURE
%% ============================================================================

\chapter{Noetic Curvature}
\label{ch:noetic-curvature}

\begin{tcolorbox}[colback=blue!5!white,colframe=blue!75!black,title=\vnewfour]
This chapter develops the theory of curvature for noetic manifolds, measuring
how parallel transport around closed loops fails to return vectors to
their original position.
\end{tcolorbox}

\section{The Riemann Curvature Tensor}

\begin{definition}[Noetic Riemann Curvature]
\label{def:riemann-curvature}
The \textbf{noetic Riemann curvature tensor} $\NoeticCurvature$ is defined by:
\[
    \NoeticCurvature(X, Y)Z := \NoeticConnection_X \NoeticConnection_Y Z
    - \NoeticConnection_Y \NoeticConnection_X Z - \NoeticConnection_{[X,Y]} Z
\]
This measures the non-commutativity of covariant differentiation.
\end{definition}

\begin{proposition}[Components of Curvature]
\label{prop:curvature-components}
In local coordinates:
\[
    R^\nu{}^l{}_{ijk} = \frac{\partial \Gamma^\nu{}^l_{jk}}{\partial x^i}
    - \frac{\partial \Gamma^\nu{}^l_{ik}}{\partial x^j}
    + \Gamma^\nu{}^l_{im} \Gamma^\nu{}^m_{jk}
    - \Gamma^\nu{}^l_{jm} \Gamma^\nu{}^m_{ik}
\]
\end{proposition}

\begin{theorem}[Curvature Symmetries]
\label{thm:curvature-symmetries}
The noetic Riemann curvature satisfies:
\begin{enumerate}[label=(\alph*)]
    \item Antisymmetry: $R^\nu{}^l{}_{ijk} = -R^\nu{}^l{}_{jik}$
    \item First Bianchi: $R^\nu{}^l{}_{ijk} + R^\nu{}^l{}_{jki} + R^\nu{}^l{}_{kij} = 0$
    \item With metric: $R^\nu_{ijkl} = -R^\nu_{ijlk} = R^\nu_{klij}$
\end{enumerate}
\end{theorem}

\section{Ricci and Scalar Curvature}

\begin{definition}[Noetic Ricci Curvature]
\label{def:ricci-curvature}
The \textbf{noetic Ricci tensor} is the trace:
\[
    \mathrm{Ric}^\nu_{ij} := R^\nu{}^k{}_{ikj}
\]
\end{definition}

\begin{definition}[Noetic Scalar Curvature]
\label{def:scalar-curvature}
The \textbf{noetic scalar curvature} is:
\[
    R^\nu := g^{\nu ij} \mathrm{Ric}^\nu_{ij}
\]
\end{definition}

\begin{theorem}[Contracted Bianchi Identity]
\label{thm:contracted-bianchi}
The Einstein tensor $G^\nu_{ij} := \mathrm{Ric}^\nu_{ij} - \frac{1}{2}R^\nu g^\nu_{ij}$
is divergence-free:
\[
    \NoeticConnection^j G^\nu_{ij} = 0
\]
\end{theorem}

\section{Sectional and Gaussian Curvature}

\begin{definition}[Sectional Curvature]
\label{def:sectional-curvature}
For linearly independent $u, v \in T_pM$, the \textbf{sectional curvature} is:
\[
    K^\nu(u, v) := \frac{\langle \NoeticCurvature(u, v)v, u \rangle_\nu}
    {\|u\|_\nu^2 \|v\|_\nu^2 - \langle u, v \rangle_\nu^2}
\]
\end{definition}

\begin{proposition}[Curvature Interpretation]
\label{prop:curvature-interpretation}
The sectional curvature $K^\nu(u, v)$ measures:
\begin{itemize}
    \item $K^\nu > 0$: Positive curvature (geodesics converge)
    \item $K^\nu = 0$: Flat (Euclidean behavior locally)
    \item $K^\nu < 0$: Negative curvature (geodesics diverge)
\end{itemize}
In noetic terms, positive curvature indicates ideas that ``attract,''
negative curvature indicates ideas that ``repel.''
\end{proposition}

\begin{theorem}[Gauss-Bonnet for Noetic Surfaces]
\label{thm:gauss-bonnet}
For a compact noetic 2-manifold $\NoeticManifold$ with Gaussian curvature $K^\nu$:
\[
    \int_M K^\nu \, dA + \int_{\partial M} \kappa_g \, ds = 2\pi \chi(M)
\]
where $\kappa_g$ is the geodesic curvature of the boundary and $\chi(M)$
is the Euler characteristic.
\end{theorem}

%% ============================================================================
%% CHAPTER 5: NOETIC GEODESICS
%% ============================================================================

\chapter{Noetic Geodesics}
\label{ch:noetic-geodesics}

\begin{tcolorbox}[colback=blue!5!white,colframe=blue!75!black,title=\vnewfour]
This chapter develops geodesics as curves of shortest distance in noetic
manifolds, representing the most ``direct'' paths between ideas.
\end{tcolorbox}

\section{Geodesic Equation}

\begin{definition}[Noetic Geodesic]
\label{def:noetic-geodesic}
A curve $\NoeticGeodesic : [a, b] \to \NoeticManifold$ is a \textbf{noetic geodesic}
if it parallel transports its own tangent vector:
\[
    \NoeticConnection_{\dot{\gamma}_\nu} \dot{\gamma}_\nu = 0
\]
\end{definition}

\begin{proposition}[Geodesic Equation in Coordinates]
\label{prop:geodesic-eq}
In local coordinates, the geodesic equation is:
\[
    \frac{d^2 \gamma^k}{dt^2} + \Gamma^\nu{}^k_{ij}
    \frac{d\gamma^i}{dt} \frac{d\gamma^j}{dt} = 0
\]
This is a system of second-order ODEs.
\end{proposition}

\begin{theorem}[Existence and Uniqueness]
\label{thm:geodesic-existence}
For any $p \in \NoeticManifold$ and $v \in T_pM$, there exists a unique
maximal geodesic $\NoeticGeodesic : I \to M$ with $\NoeticGeodesic(0) = p$
and $\dot{\gamma}_\nu(0) = v$.
\end{theorem}

\section{Exponential Map}

\begin{definition}[Exponential Map]
\label{def:exp-map}
The \textbf{exponential map} $\exp_p^\nu : T_pM \to M$ is defined by:
\[
    \exp_p^\nu(v) := \NoeticGeodesic_v(1)
\]
where $\NoeticGeodesic_v$ is the geodesic with $\NoeticGeodesic_v(0) = p$
and $\dot{\gamma}_\nu(0) = v$.
\end{definition}

\begin{theorem}[Local Diffeomorphism]
\label{thm:exp-local-diffeo}
The exponential map $\exp_p^\nu$ is a local diffeomorphism near $0 \in T_pM$.
\end{theorem}

\begin{definition}[Normal Coordinates]
\label{def:normal-coords}
\textbf{Normal coordinates} at $p$ are coordinates induced by $\exp_p^\nu$
through a basis of $T_pM$. In these coordinates:
\begin{itemize}
    \item $\Gamma^\nu{}^k_{ij}(p) = 0$
    \item Geodesics through $p$ are straight lines
    \item $g^\nu_{ij}(p) = \delta_{ij}$
\end{itemize}
\end{definition}

\section{Geodesic Distance and Completeness}

\begin{definition}[Injectivity Radius]
\label{def:injectivity-radius}
The \textbf{injectivity radius} at $p$ is:
\[
    \mathrm{inj}(p) := \sup\{r > 0 : \exp_p^\nu|_{B_r(0)} \text{ is injective}\}
\]
\end{definition}

\begin{definition}[Geodesic Completeness]
\label{def:geodesic-complete}
$\NoeticManifold$ is \textbf{geodesically complete} if every geodesic
can be extended to all of $\mathbb{R}$.
\end{definition}

\begin{theorem}[Hopf-Rinow for Noetic Manifolds]
\label{thm:hopf-rinow}
For a connected noetic manifold, the following are equivalent:
\begin{enumerate}[label=(\roman*)]
    \item $\NoeticManifold$ is geodesically complete
    \item $(\NoeticManifold, d_\nu)$ is a complete metric space
    \item Closed bounded sets are compact
    \item Any two points can be joined by a minimal geodesic
\end{enumerate}
\end{theorem}

\section{Jacobi Fields and Conjugate Points}

\begin{definition}[Jacobi Field]
\label{def:jacobi-field}
A \textbf{Jacobi field} $J$ along a geodesic $\NoeticGeodesic$ is a vector
field satisfying the Jacobi equation:
\[
    \NoeticConnection_{\dot{\gamma}} \NoeticConnection_{\dot{\gamma}} J
    + \NoeticCurvature(J, \dot{\gamma})\dot{\gamma} = 0
\]
\end{definition}

\begin{definition}[Conjugate Point]
\label{def:conjugate-point}
Points $p = \NoeticGeodesic(t_0)$ and $q = \NoeticGeodesic(t_1)$ are
\textbf{conjugate} along $\NoeticGeodesic$ if there exists a non-zero
Jacobi field $J$ with $J(t_0) = J(t_1) = 0$.
\end{definition}

\begin{theorem}[Conjugate Points and Minimality]
\label{thm:conjugate-minimal}
A geodesic $\NoeticGeodesic$ fails to minimize distance beyond its first
conjugate point.
\end{theorem}

%% ============================================================================
%% CHAPTER 6: PARALLEL TRANSPORT AND HOLONOMY
%% ============================================================================

\chapter{Parallel Transport and Holonomy}
\label{ch:holonomy}

\begin{tcolorbox}[colback=blue!5!white,colframe=blue!75!black,title=\vnewfour]
This chapter develops parallel transport along curves and the holonomy
group, which measures the global non-triviality of the connection.
\end{tcolorbox}

\section{Parallel Transport}

\begin{definition}[Parallel Vector Field]
\label{def:parallel-field}
A vector field $V$ along a curve $\gamma$ is \textbf{parallel} if:
\[
    \NoeticConnection_{\dot{\gamma}} V = 0
\]
\end{definition}

\begin{theorem}[Existence of Parallel Transport]
\label{thm:parallel-transport-exists}
For any curve $\gamma : [a, b] \to M$ and $v \in T_{\gamma(a)}M$, there
exists a unique parallel vector field $V$ along $\gamma$ with $V(a) = v$.
\end{theorem}

\begin{definition}[Parallel Transport Map]
\label{def:parallel-transport-map}
The \textbf{parallel transport} along $\gamma$ from $\gamma(a)$ to $\gamma(b)$ is:
\[
    \ParallelTransport : T_{\gamma(a)}M \to T_{\gamma(b)}M, \quad v \mapsto V(b)
\]
where $V$ is the parallel extension of $v$.
\end{definition}

\begin{proposition}[Properties of Parallel Transport]
\label{prop:parallel-properties}
Parallel transport satisfies:
\begin{enumerate}[label=(\alph*)]
    \item Linearity: $\ParallelTransport(av + bw) = a\ParallelTransport(v) + b\ParallelTransport(w)$
    \item Isometry: $\langle \ParallelTransport v, \ParallelTransport w \rangle_\nu
          = \langle v, w \rangle_\nu$ (for metric-compatible connections)
    \item Functoriality: $\mathcal{P}_{\gamma_2} \circ \mathcal{P}_{\gamma_1}
          = \mathcal{P}_{\gamma_2 \cdot \gamma_1}$ (for concatenated curves)
\end{enumerate}
\end{proposition}

\section{Holonomy Groups}

\begin{definition}[Holonomy Group]
\label{def:holonomy-group}
The \textbf{holonomy group} at $p \in M$ is:
\[
    \HolonomyGroup_p := \{\ParallelTransport : \gamma \text{ is a loop based at } p\}
    \subseteq \mathrm{GL}(T_pM)
\]
For metric-compatible connections: $\HolonomyGroup_p \subseteq O(T_pM)$.
\end{definition}

\begin{definition}[Restricted Holonomy]
\label{def:restricted-holonomy}
The \textbf{restricted holonomy group} $\mathrm{Hol}^0(\NoeticConnection)_p$
consists of parallel transports around null-homotopic loops.
\end{definition}

\begin{theorem}[Holonomy and Curvature]
\label{thm:holonomy-curvature}
$\mathrm{Hol}^0(\NoeticConnection) = \{e\}$ if and only if $\NoeticCurvature = 0$.
That is, the connection is flat iff the restricted holonomy is trivial.
\end{theorem}

\begin{proposition}[Ambrose-Singer]
\label{prop:ambrose-singer}
The Lie algebra of $\HolonomyGroup$ is spanned by elements of the form
$\ParallelTransport^{-1} \NoeticCurvature(X, Y) \ParallelTransport$ for all
$X, Y$ tangent vectors and all parallel transports $\ParallelTransport$.
\end{proposition}

\section{Flat Connections and Global Parallelism}

\begin{definition}[Flat Connection]
\label{def:flat-connection}
A connection is \textbf{flat} if $\NoeticCurvature = 0$.
\end{definition}

\begin{theorem}[Global Parallelism]
\label{thm:global-parallelism}
On a simply connected manifold $\NoeticManifold$ with flat connection:
\begin{enumerate}[label=(\roman*)]
    \item Parallel transport is path-independent
    \item There exists a global parallel frame
    \item The manifold is locally isometric to $\mathbb{R}^n$ with the flat metric
\end{enumerate}
\end{theorem}

\begin{example}[Noetic Torus]
\label{ex:flat-torus}
The noetic torus $\mathbb{T}^n_\nu = \mathbb{R}^n/\mathbb{Z}^n$ with the flat
metric inherited from $\mathbb{R}^n$ has:
\begin{itemize}
    \item $\NoeticCurvature = 0$ (flat)
    \item $\HolonomyGroup = \{e\}$ (trivial holonomy)
    \item Non-trivial fundamental group $\pi_1 \cong \mathbb{Z}^n$
\end{itemize}
\end{example}

%% ============================================================================
%% CHAPTER 7: APPLICATIONS TO TKS
%% ============================================================================

\chapter{Applications to TKS}
\label{ch:geometry-tks-applications}

\begin{tcolorbox}[colback=green!5!white,colframe=green!75!black,title=Synthesis]
This chapter synthesizes differential noetic geometry with the broader TKS
framework, showing how geometric concepts illuminate noetic structures.
\end{tcolorbox}

\section{Idea Space as Noetic Manifold}

\begin{construction}[Geometric Idea Space]
\label{constr:geometric-idea-space}
The idea space $\Ideas$ carries a natural noetic manifold structure:
\begin{enumerate}[label=(\roman*)]
    \item Local coordinates: Components of ideas in a basis
    \item Noetic metric: Similarity measure between ideas
    \item Connection: How ideas ``relate'' infinitesimally
\end{enumerate}
\end{construction}

\begin{proposition}[Noetic Distance Interpretation]
\label{prop:noetic-distance-interpretation}
The noetic distance $d_\nu(\iota, \iota')$ measures:
\begin{itemize}
    \item Conceptual dissimilarity between ideas
    \item Effort required to transform one idea into another
    \item Geodesic paths represent ``natural'' transitions
\end{itemize}
\end{proposition}

\section{Geodesics as Conceptual Paths}

\begin{theorem}[Geodesics of Ideation]
\label{thm:geodesics-ideation}
A noetic geodesic $\NoeticGeodesic$ represents:
\begin{enumerate}[label=(\alph*)]
    \item The shortest conceptual path between ideas
    \item A path of minimal ``cognitive effort''
    \item Natural interpolation in idea space
\end{enumerate}
\end{theorem}

\begin{example}[Analogical Reasoning]
\label{ex:analogical-reasoning}
Given ideas $\iota_A, \iota_B$ with analogy $\iota_A : \iota_{A'} :: \iota_B : ?$,
the answer $\iota_{B'}$ lies on the geodesic:
\[
    \iota_{B'} = \exp_{\iota_B}^\nu\left(\ParallelTransport_{\iota_A \to \iota_B}
    (\log_{\iota_A}^\nu(\iota_{A'}))\right)
\]
where $\log^\nu$ is the inverse of $\exp^\nu$.
\end{example}

\section{Curvature and Conceptual Structure}

\begin{proposition}[Curvature Interpretation]
\label{prop:curvature-conceptual}
The curvature of idea space at $\iota$ reflects:
\begin{itemize}
    \item $K^\nu > 0$: Ideas in this region ``cluster'' (related concepts attract)
    \item $K^\nu < 0$: Ideas ``diverge'' (conceptual space is sparse)
    \item $K^\nu = 0$: Linear, Euclidean-like conceptual structure
\end{itemize}
\end{proposition}

\begin{definition}[Conceptual Focusing]
\label{def:conceptual-focusing}
In regions of positive curvature, geodesics from an idea $\iota$ converge,
leading to \textbf{conceptual focusing}: different paths of thought lead
to similar conclusions.
\end{definition}

\begin{definition}[Conceptual Divergence]
\label{def:conceptual-divergence}
In regions of negative curvature, geodesics diverge exponentially, representing
\textbf{conceptual divergence}: small differences in initial direction lead
to vastly different ideas.
\end{definition}

\section{Holonomy and Circular Reasoning}

\begin{theorem}[Holonomy Interpretation]
\label{thm:holonomy-interpretation}
Non-trivial holonomy around a loop in idea space represents:
\begin{enumerate}[label=(\roman*)]
    \item Circular reasoning that returns with changed perspective
    \item The Hermeneutic circle: understanding changes through cyclic interpretation
    \item Path-dependence in conceptual development
\end{enumerate}
\end{theorem}

\begin{example}[Hermeneutic Loop]
\label{ex:hermeneutic}
Following a loop $\gamma$ through ideas $\iota_1 \to \iota_2 \to \cdots \to \iota_n \to \iota_1$,
the parallel transport $\ParallelTransport \neq \mathrm{id}$ represents how
our understanding of $\iota_1$ is transformed by the journey through other ideas.
\end{example}

\section{Integration with v7.3 Structures}

\begin{tcolorbox}[colback=yellow!5!white,colframe=yellow!75!black,title=Cross-Reference: v7.3]
\textbf{Integration with Transfinite Fractals:} The noetic manifold structure
is compatible with the transfinite fractal hierarchy. Each fractal level
$\TransFrac{\alpha}$ inherits a noetic metric, and the refinement maps
preserve geodesic structure locally.
\end{tcolorbox}

\begin{proposition}[Fractal-Geometric Compatibility]
\label{prop:fractal-geometric}
For a transfinite fractal $\{\TransFrac{\alpha}\}_{\alpha < \kappa}$:
\begin{enumerate}[label=(\roman*)]
    \item Each $\TransFrac{\alpha}$ is a noetic sub-manifold
    \item Inclusion maps $\TransFrac{\alpha} \hookrightarrow \TransFrac{\beta}$ are isometric embeddings
    \item The colimit $\TransFrac{\kappa}$ carries a limiting noetic metric
\end{enumerate}
\end{proposition}

\begin{tcolorbox}[colback=yellow!5!white,colframe=yellow!75!black,title=Cross-Reference: Higher Categories]
\textbf{Integration with $(\infty,1)$-Categories:} The noetic quasi-category
$\mathcal{N}_\infty$ has a geometric enrichment where mapping spaces carry
noetic metrics making $\mathcal{N}_\infty$ into a \textbf{geodesic $\infty$-category}.
\end{tcolorbox}

\begin{definition}[Geodesic $\infty$-Category]
\label{def:geodesic-infty-cat}
A quasi-category $\mathcal{C}$ is \textbf{geodesically enriched} if:
\begin{enumerate}[label=(\roman*)]
    \item Each mapping space $\mathrm{Map}(x, y)$ is a noetic manifold
    \item Composition is compatible with geodesic structure
    \item Higher simplices are geodesically interpolated
\end{enumerate}
\end{definition}

\begin{tcolorbox}[colback=yellow!5!white,colframe=yellow!75!black,title=Cross-Reference: Domain Theory]
\textbf{Integration with Domain Theory:} Scott domains can be equipped with
a generalized noetic metric where:
\[
    d_\nu(x, y) := \inf\{\epsilon > 0 : x \sqsubseteq_\epsilon y \land y \sqsubseteq_\epsilon x\}
\]
This makes $(D, d_\nu)$ a quasi-metric space compatible with the Scott topology.
\end{tcolorbox}

\section{Noetic Einstein Equations}

\begin{definition}[Noetic Stress-Energy]
\label{def:noetic-stress-energy}
The \textbf{noetic stress-energy tensor} $T^\nu_{ij}$ encodes the distribution
of ``conceptual matter'' in idea space:
\begin{itemize}
    \item $T^\nu_{00}$: Density of ideas (conceptual mass)
    \item $T^\nu_{0i}$: Flow of ideas (conceptual momentum)
    \item $T^\nu_{ij}$: Conceptual pressure and stress
\end{itemize}
\end{definition}

\begin{axiom}[Noetic Einstein Equations]
\label{axiom:noetic-einstein}
The geometry of idea space is determined by its conceptual content:
\[
    G^\nu_{ij} + \Lambda^\nu g^\nu_{ij} = 8\pi G_\nu T^\nu_{ij}
\]
where $G^\nu$ is the noetic Einstein tensor, $\Lambda^\nu$ is the noetic
cosmological constant, and $G_\nu$ is the noetic gravitational constant.
\end{axiom}

\begin{theorem}[Conceptual Dynamics]
\label{thm:conceptual-dynamics}
The noetic Einstein equations imply:
\begin{enumerate}[label=(\roman*)]
    \item Ideas curve the surrounding conceptual space
    \item Geodesics (natural thought paths) bend toward ``heavy'' ideas
    \item Conceptual ``black holes'' may exist: ideas of such density that
          no geodesic escapes
\end{enumerate}
\end{theorem}

%% ============================================================================
%% END OF GEOMETRY TRACK
%% ============================================================================

\begin{tcolorbox}[colback=blue!5!white,colframe=blue!75!black,title=Track Summary,breakable]
The Differential Noetic Geometry track has established:
\begin{enumerate}
    \item Noetic manifolds as smooth spaces of ideas
    \item The noetic metric measuring conceptual distance
    \item Connections for comparing ideas at different points
    \item Curvature measuring local conceptual structure
    \item Geodesics as optimal conceptual paths
    \item Holonomy capturing the effects of circular reasoning
    \item Applications connecting geometry to TKS semantics
\end{enumerate}
This provides a complete differential-geometric foundation for
analyzing the structure and dynamics of idea space.
\end{tcolorbox}


%% ============================================================================
%% PART IV: CORRECTED CHAPTERS
%% ============================================================================
\part{Corrected Chapters}
\label{part:corrected}

\input{chapter27_corrected.tex}

\input{geometry_corrected.tex}

\end{document}
